\section{Basic}

\begin{definition}
  An \textit{(undirected) graph} is a pair $G = (V,E)$ consisting of a set $V$ 
  and a multi-set $E$ of unordered pairs of elements in $V$. 
  Elements of $V$ are called \textit{vertices} and elements of $E$ are called \textit{edges}.
  For two vertices $u,v \in V$, we say $u$ and $v$ are \textit{adjacent} if $\{u,v\} \in E$.
  We call an undirected graph $G$ \textit{simple} if all edges are distinct (i.e., $E$ is a set), 
  and each edge is a pair of distinct vertices.
\end{definition}

\begin{definition}
  A \textit{directed graph} is a pair $G = (V,E)$ consisting of a set $V$ 
  and a multi-set $E$ of ordered pairs of elements in $V$. 
  Elements of $V$ are called \textit{vertices} and elements of $E$ are called \textit{edges}.
  We call a directed graph $G$ \textit{simple} if all edges are distinct and
  each edge is a pair of distinct vertices.
\end{definition}
